\documentclass[11pt]{article}
\usepackage{amsmath,hyperref}
\title{A current rival report can miss the blocker after a change}
\author{}
\date{}
\begin{document}
\maketitle

A rail signal can report the same rival in two states and give different decisions after that rival is withdrawn. One state admits the original train; in the other a second rival still blocks it. Removing an approach can even turn an admission into a refusal by exposing a higher-priority candidate from behind the old head of its link. These are executions of selected SUMO methods with supplied approaches and route predicates, not results for reachable full-network configurations.

\paragraph{What the source reports.}
The C++ harness extracts exact methods from \href{https://github.com/eclipse-sumo/sumo/tree/fa7df22d7f1c4ae245192ecc32410a49f1ffb2a3}{SUMO v1.26.0}: driveway reservation, link-conflict checking, yielding priority, closest approach and the rail diagnostic path. Geometry, approach lists, occupancy and overlap are supplied by the fixture; it does not simulate train motion or claim that these approach states were observed in a network. The code first checks occupancy, then runs the competing links in order, stopping at the first required yield. Each link supplies its closest approach. The diagnostic reruns this path, resetting stored vectors; it can write the active train when reservation succeeds. The native reserve call is made separately from the diagnostic call.

In one executed pair the vector shows P and omits Q, even when Q is present on a later link. Deleting P clears one state but leaves Q as a blocker in the other. In another test, a nonblocking approach leads a more urgent approach on the same link; the diagnostic initially has no priority rival. Removing the nearer head promotes the other approach, and the gate refuses the reservation. A separate occupancy block yields no rival entries and also refuses. Thus this diagnostic has neither a completeness guarantee for alternative futures nor an effect that is interchangeable with physical clearance.

\paragraph{Cutoff for bounded removals.}
Fix the focal claim, a clear occupancy gate, ordered queues at each link, and one immutable eligibility bit per candidate: if selected as head, is it a conflict requiring the focal train to yield? The admissible revision removes at most $k$ named candidates altogether. It does not add approaches, reorder queues, or alter any survivor's physical, route or priority predicate. With these premises the call is blocked if any queue's first surviving candidate has bit one. For exact queries it suffices to retain the first $k+1$ candidates and their bits at each link, or the full shorter queue. Deeper candidates cannot be exposed by the allowed deletions. On an isolated link with at least $k+1$ positions, exact answers on deletions of prefixes determine every one of these bits, so a summary uniformly covering all assignments needs at least $k+1$ bits besides the fixed identities. No additive lower bound over several links is asserted.

For the separate question whether \emph{some} allowed deletion admits the claim, let $b_l$ be the length of the initial consecutive blocking prefix on link $l$. Any admitting deletion must remove every such candidate, since leaving one makes the earliest survivor blocking. Removing just the specified prefixes suffices: each remaining head is nonblocking or the queue is empty. Thus the exact minimum deletion count is $\sum_l b_l$. Any true prefixes totaling more than $k$ entries certify obstruction under the declared budget, without giving the whole queue. For a successful candidate, retaining only its completed predecessor or the current reason for refusal is insufficient if the removed head can reveal another candidate. The exact query and minimum-deletion summary answer different questions.

\paragraph{Checks and scope.}
The selected C++ methods verify the equal-diagnostic, hidden lower-link rival, same-link promotion and empty-report occupancy controls. A separate exact Python enumeration checks the fixed-queue formula on 83,952 deletion queries and 3,600 budget instances against full independent scans, with 5,436 queries turning an initial admission into a refusal. The strongest same-input direct approach scans the static queue heads or maintains them with an index and ties. The new information is which negatives and queue positions a reusable guard must cover, not an algorithmic advantage over that reference. The bounded law does not certify how route state and physical continuity change after a real release or reroute; those can alter both the queue and the eligibility bits. An older rival report without complete replacement evidence cannot justify advance admission. This does not discharge the independent source or architecture obligations.
\end{document}
